\clearpage
\section{Design}

The middleware follows a \textbf{token-based design}. Each logical critical
section is associated with one durable RabbitMQ queue, named
\texttt{cs\_token\_<csName>}. When the critical section is free, that queue
contains exactly one persistent token message. When a process enters the
section, it consumes that token with manual acknowledgements and keeps the
delivery unacknowledged for the whole duration of the critical section.

The public API exposes three operations: \texttt{enter()}, \texttt{exit()}, and
\texttt{close()}. Client code does not interact directly with queues,
consumers, acknowledgements, or failure recovery; these concerns are managed by
the middleware implementation.

\subsection{Bootstrap and ownership}

The only delicate phase is token creation. To prevent concurrent starters from
creating duplicate tokens, the middleware uses a second temporary exclusive
queue, \texttt{cs\_token\_creation\_guard\_<csName>}. This queue is not the
critical section itself; it is only a broker-side guard that serializes the
decision ``does the first token need to be created?''

Once bootstrap is complete, the token queue becomes the only relevant shared
state. If the token is present in the queue, the section is free. If the token
has been delivered and is still unacknowledged, some process owns the section.
For the same reason, a clean shutdown while holding the token is treated as a
protected release transition: the middleware serializes that release with the
same broker-side guard used during bootstrap, so a concurrent initializer
cannot observe an inconsistent intermediate state.

\subsection{Failure recovery}

The recovery strategy is broker-driven. If the current owner crashes or loses
its channel before calling \texttt{exit()}, RabbitMQ automatically requeues the
unacknowledged token. The middleware therefore avoids a second recovery
protocol: crash tolerance is obtained by aligning the critical-section protocol
with standard RabbitMQ delivery semantics.
